\documentclass[12pt]{article}
\usepackage[T1]{fontenc}
\usepackage[utf8]{inputenc}
\usepackage{lmodern}
\usepackage{textcomp}
\usepackage{enumitem}
\usepackage{geometry}
\geometry{margin=1in}
\begin{document}
\begin{center}
Answer Key - Biology - Undergraduate Level Assessment 5 \
\small (Answers are ordered exactly as in the question paper.)
\end{center}

\section*{Multiple Choice}
\begin{enumerate}[label=\arabic*.]
\item \textbf{Answer:} C
\\ \textit{Options:} A) Lipids are the primary source of energy for cellular processes.; B) Lipids are soluble in water but insoluble in organic solvents.; C) Phospholipids are water soluble.; D) Steroids are lipids that consist of glycerol and fatty acids.; E) Lipids make up most cell surface receptors.
\\ \textit{Note:} Phospholipids have a hydrophilic (water-attracting) phosphate head and two hydrophobic (water-repelling) fatty acid tails, allowing them to form bilayers that are soluble in water, thereby making them essential components of cell membranes.
\item \textbf{Answer:} A
\\ \textit{Options:} A) a universally accepted fact about the origin and development of all species.; B) a hypothesis in the process of being tested and verified.; C) a disproven theory about the development of species over time.; D) a speculation about possible changes in populations over time.; E) an opinion that some scientists hold about how living things change over time.
\\ \textit{Note:} The theory of evolution is supported by extensive scientific evidence from various fields, including genetics, paleontology, and comparative anatomy, establishing it as a foundational principle that explains the diversity of life and the common descent of all species.
\item \textbf{Answer:} D
\\ \textit{Options:} A) By observing the cell under a microscope; B) Using fluorescent tagging of proteins; C) Tracking changes in cell size and shape; D) Autoradiography with labelled amino acids; E) Through protein synthesis
\\ \textit{Note:} Autoradiography with labeled amino acids allows researchers to trace the incorporation of specific amino acids into newly synthesized proteins, thereby demonstrating the role of amino acids in protein construction while also showing the degradation of existing proteins in catabolic processes.
\item \textbf{Answer:} A
\\ \textit{Options:} A) Plants, animals, and fungi all have both chloroplasts and mitochondria.; B) Plants and fungi have chloroplasts but no mitochondria; animals have only mitochondria.; C) Animals have both chloroplasts and mitochondria; plants and fungi have only chloroplasts.; D) Plants and animals have chloroplasts but no mitochondria; fungi have only mitochondria.; E) Plants and animals have mitochondria but no chloroplasts; fungi have only chloroplasts.
\\ \textit{Note:} The correct answer is misleading because while plants possess both chloroplasts and mitochondria, animals and fungi only have mitochondria, thus making the statement inaccurate as a generalization.
\item \textbf{Answer:} C
\\ \textit{Options:} A) Appetitive behavior, exploratory behavior, quiescence; B) Termination, appetitive behavior, exploratory behavior; C) Appetitive behavior, consummatory act, quiescence; D) Consummatory act, exploratory behavior, termination; E) Appetitive behavior, termination, consummatory act
\\ \textit{Note:} The correct answer identifies the typical phases of a behavioral act, where appetitive behavior involves seeking or approaching a goal, the consummatory act is the actual performance of the behavior, and quiescence refers to the resting state following the behavior, reflecting the natural progression of behavior in biological contexts.
\end{enumerate}

\section*{True / False}
\begin{enumerate}[label=\arabic*.]
\setcounter{enumi}{5}
\item \textbf{Answer:} False
\\ \textit{Options:} A) True; B) False
\\ \textit{Note:} Using negative reinforcement alone is insufficient to assess a dog's hearing sensitivity because it does not adequately measure the dog's ability to detect and differentiate between various sound frequencies, which requires a more comprehensive approach that includes positive reinforcement and controlled testing conditions.
\item \textbf{Answer:} False
\\ \textit{Options:} A) True; B) False
\\ \textit{Note:} Down's syndrome is caused by the presence of an extra copy of chromosome 21, not a mutation in chromosome 14.
\item \textbf{Answer:} True
\\ \textit{Options:} A) True; B) False
\\ \textit{Note:} Tracheids and vessel elements are specialized cells in the xylem that facilitate the transport of water and dissolved minerals while also contributing to the structural integrity of the plant.
\item \textbf{Answer:} False
\\ \textit{Options:} A) True; B) False
\\ \textit{Note:} Stabilizing selection favors intermediate traits and reduces variability in reproductive success by selecting against extreme phenotypes, thus promoting a more uniform population.
\item \textbf{Answer:} False
\\ \textit{Options:} A) True; B) False
\\ \textit{Note:} The rise in intracellular free calcium in the sea urchin oocyte initiates the acrosomal reaction, but this process is crucial for the fertilization of the egg by a single sperm, not multiple sperm.
\end{enumerate}

\section*{Long Form Response}
\begin{enumerate}[label=\arabic*.]
\setcounter{enumi}{10}
\item \textbf{Answer:} Courtship behaviors, such as the male common tern presenting a fish, play a crucial role in communication and reproductive success among animal species. These behaviors serve as a means of sexual identification, allowing potential mates to recognize their species and assess the fitness of partners. By offering a fish, the male not only signals his ability to provide resources but also demonstrates his health and vitality, which are attractive traits for females seeking to maximize reproductive success. In many cases, these displays are essential for initiating mating behaviors and can influence female choice, ultimately impacting egg production and the continuation of species. Thus, courtship behaviors are fundamental in facilitating reproductive opportunities and ensuring genetic diversity within populations.
\\ \textit{Note:} correct because it emphasizes how courtship behaviors, such as the male common tern presenting a fish, facilitate species recognition and communicate the male's fitness and resource availability, ultimately enhancing reproductive success for both partners.
\item \textbf{Answer:} Living organisms are distinguished from non-living matter by their complex organization and their ability to perform essential life functions. Key characteristics include the structured arrangement of cells, which are the basic units of life, and the capacity to consume energy for growth and maintenance. Living things exhibit movement, whether voluntary or involuntary, and respond to environmental stimuli, indicating a level of interaction with their surroundings. Furthermore, they possess the ability to grow, reproduce, and adapt to changing conditions, ensuring survival and evolution over time. These features underscore the intricate functionality required for life, setting living organisms apart from inanimate objects.
\\ \textit{Note:} correct because it identifies the complex cellular organization, energy consumption, movement, and responsiveness to stimuli as fundamental characteristics that differentiate living organisms from non-living matter, highlighting their essential life functions.
\end{enumerate}

\vfill
\noindent\textit{End of Answer Key}
\end{document}